\chapter{La symétrie}\label{ch:g4:symmetry}

Papillons, visages, flocons de neige~: la nature aime les figures
dont les deux moitiés coïncident quand on les plie. Ce chapitre
explore la symétrie avec des plis et du papier quadrillé --- les
constructions à l'équerre attendent \cref{ch:g6:symmetry}.

\section{Plier et faire coïncider}

\begin{definition}[Axe de symétrie]\label{def:g4:symmetry:def}
Une droite est un \emph{axe de symétrie}\index{axe de symétrie}
d'une figure lorsque le pliage de la figure le long de cette droite
fait coïncider les deux moitiés \emph{exactement}. Une figure peut
n'avoir aucun axe, un seul, ou plusieurs.
\end{definition}

\begin{omfigure}
\begin{tikzpicture}[scale=0.6]
  % heart-ish symmetric shape: isosceles triangle
  \draw[thick, fill=omDef!12] (0,0) -- (2.4,0) -- (1.2,2.6) -- cycle;
  \draw[omThm, dashed, thick] (1.2,-0.5) -- (1.2,3.0);
  \node[below, font=\small] at (1.2,-0.7) {un axe};
  % rectangle: two axes
  \begin{scope}[xshift=5.4cm]
  \draw[thick, fill=omDef!12] (0,0.3) rectangle (3.2,2.3);
  \draw[omThm, dashed, thick] (1.6,-0.4) -- (1.6,3.0);
  \draw[omThm, dashed, thick] (-0.5,1.3) -- (3.7,1.3);
  \node[below, font=\small] at (1.6,-0.7) {deux axes};
  \end{scope}
  % scalene: none
  \begin{scope}[xshift=11.2cm]
  \draw[thick, fill=omDef!12] (0,0) -- (3.0,0.4) -- (0.8,2.2) -- cycle;
  \node[below, font=\small] at (1.4,-0.7) {aucun axe};
  \end{scope}
\end{tikzpicture}

{\small Plier le long du trait en pointillés~: les moitiés
coïncident. Le dernier triangle n'a aucun pli qui fonctionne ---
pas d'axe de symétrie.}
\end{omfigure}

\begin{example}[Tester avec du papier calque]\label{ex:g4:symmetry:trace}
Calquer la figure, retourner le calque le long de l'axe candidat, et
le reposer~: si le calque recouvre exactement la figure, l'axe est
authentique. L'œil seul se trompe souvent --- la diagonale d'un
rectangle \emph{semble} un axe, mais le test du pli dit non
(essayez~!).
\end{example}

\section{Compléter des figures sur un quadrillage}

\begin{method}[Symétrique d'une figure par rapport à une ligne du
quadrillage]\label{met:g4:symmetry:grid}
L'axe est une ligne verticale (ou horizontale) du quadrillage.
\begin{enumerate}
  \item Prendre chaque sommet de la figure à tour de rôle~;
  \item compter sa distance à l'axe en carreaux~;
  \item placer le sommet image à la \emph{même} distance de
        \emph{l'autre} côté, sur la même ligne (ou colonne)~;
  \item joindre les sommets images dans le même ordre, et
        vérifier~: l'image est le jumeau miroir de la figure.
\end{enumerate}
\end{method}

\begin{omfigure}
\begin{tikzpicture}[scale=0.55]
  \draw[gray!40, thin] (0,0) grid (12,5);
  \draw[omThm, very thick] (6,-0.4) -- (6,5.4);
  \draw[thick, omDef, fill=omDef!15]
    (1,1) -- (5,1) -- (5,2) -- (3,2) -- (3,4) -- (1,4) -- cycle;
  \draw[thick, omMeth, fill=omMeth!15]
    (11,1) -- (7,1) -- (7,2) -- (9,2) -- (9,4) -- (11,4) -- cycle;
  \draw[dashed] (5,2) -- (7,2);
  \draw[dashed] (3,4) -- (9,4);
\end{tikzpicture}

{\small Compléter par rapport à l'axe rouge~: chaque sommet saute à
la même distance de l'autre côté. Les sommets touchant l'axe
resteraient en place.}
\end{omfigure}

\begin{example}[Même forme, même taille, retournée]\label{ex:g4:symmetry:properties}
L'image a exactement les mêmes longueurs et les mêmes angles que
l'original~: la symétrie copie la figure --- mais la retourne, comme
une main gauche et une main droite. Si la lettre originale est un F,
son image miroir est un F à l'envers, pas un autre F.
\end{example}

\section{Axes des formes usuelles}

\begin{example}[Compter les axes]\label{ex:g4:symmetry:shapes}
En pliant~:
\begin{itemize}
  \item un carré~: $4$ axes (deux par les milieux des côtés
        opposés, deux le long des diagonales)~;
  \item un rectangle~: $2$ axes (milieux des côtés opposés
        seulement~!)~;
  \item un triangle isocèle~: $1$~; un triangle équilatéral~: $3$~;
  \item un cercle~: toute droite passant par son centre --- plus
        d'axes qu'on ne peut en compter.
\end{itemize}
\end{example}

\begin{example}[Symétrie dans l'alphabet]\label{ex:g4:symmetry:letters}
A, M, T, U, V ont un axe vertical~; B, C, D, E en ont un horizontal~;
H, I, O, X ont les deux. Des mots faits des bonnes lettres peuvent
aussi être symétriques~: MUM a un axe vertical~; OXO marche dans les
deux sens.
\end{example}

\section{Exercices}

\begin{exercise}[$\star$]\label{exo:g4:symmetry:1}
Calquer et plier~: lesquelles de ces figures ont un axe de
symétrie --- un carré~; un rectangle (non carré)~; un triangle
scalène (tous les côtés différents)~; un cercle~?
\end{exercise}

\begin{exercise}[$\star$]\label{exo:g4:symmetry:2}
Dessiner les axes de symétrie de~: un carré~; un rectangle~; un
triangle équilatéral. Combien pour chacun~?
\end{exercise}

\begin{exercise}[$\star$]\label{exo:g4:symmetry:3}
Utiliser le test du papier calque de \cref{ex:g4:symmetry:trace}
pour montrer que la diagonale d'un rectangle $6 \times 4$
\emph{n'est pas} un axe de symétrie.
\end{exercise}

\begin{exercise}[$\star$]\label{exo:g4:symmetry:4}
Sur papier quadrillé, dessiner un axe vertical et la lettre L
(trois carreaux de haut, deux de large) à $2$ carreaux de l'axe.
Construire son image miroir.
\end{exercise}

\begin{exercise}[$\star$]\label{exo:g4:symmetry:5}
Sur papier quadrillé, dessiner un axe horizontal et un petit
drapeau (un triangle $1 \times 1$ sur un mât de $3$ carreaux)
au-dessus. Compléter la figure par symétrie sous l'axe.
\end{exercise}

\begin{exercise}[$\star$]\label{exo:g4:symmetry:6}
Une figure touche son axe de symétrie en un point $P$. Où se trouve
l'image de $P$~? Expliquer avec l'image du pliage.
\end{exercise}

\begin{exercise}[$\star$]\label{exo:g4:symmetry:7}
Classer les lettres capitales A, B, E, F, H, N, O, T~: axe
vertical~? axe horizontal~? les deux~? aucun~?
\end{exercise}

\begin{exercise}[$\star$]\label{exo:g4:symmetry:8}
Un triangle a pour côtés $5$~cm, $5$~cm et $3$~cm. A-t-il un axe de
symétrie~? Quels côtés le pli échange-t-il~?
\end{exercise}

\begin{exercise}[$\star\star$]\label{exo:g4:symmetry:9}
Sur papier quadrillé, dessiner une figure de son invention qui a
\emph{exactement deux} axes de symétrie, et tracer les deux axes.
(Indice~: penser à un rectangle, ou à une croix.)
\end{exercise}

\begin{exercise}[$\star\star$]\label{exo:g4:symmetry:10}
La moitié d'un dessin symétrique est donnée~: la moitié gauche d'une
maison (un mur carré et la moitié d'un toit triangulaire), avec un
axe vertical le long de son bord droit. Décrire ou dessiner la
maison complète. Comment les longueurs de la moitié droite complétée
se comparent-elles à celles de la moitié gauche~?
\end{exercise}

\begin{exercise}[$\star\star$]\label{exo:g4:symmetry:11}
Aya affirme~: «~tout triangle avec deux côtés égaux a un axe de
symétrie, et tout triangle avec un axe de symétrie a deux côtés
égaux.~» Tester sa double affirmation sur des dessins~: un triangle
isocèle, un équilatéral, un scalène. Le pliage la soutient-il~?
\end{exercise}
